%% ==============================================================================
%% sections/03-donnees-et-tikz.tex — Graphiques vectoriels, maths et code
%% ==============================================================================
\section{Données, Équations \& Code}

\begin{frame}{Visualisation vectorielle et formulation mathématique}
  \framesubtitle{Courbes \code{pgfplots} natives respectant exactement la charte chromatique}

  \begin{columns}[c, totalwidth=\textwidth]
    \begin{column}{0.51\textwidth}
      \begin{cadreFigure}
        \begin{tikzpicture}
          \begin{axis}[
            width=0.96\linewidth,
            height=4.0cm,
            xlabel={Itération $k$},
            ylabel={Précision (\%)},
            xmin=0, xmax=20,
            ymin=55, ymax=102,
            grid=major,
            grid style={dashed, ThesisMuted!25},
            legend pos=south east,
            legend style={font=\scriptsize, draw=ThesisMuted!40, fill=white},
            label style={font=\scriptsize\sffamily, text=ThesisText},
            tick label style={font=\scriptsize\sffamily, text=ThesisMuted}
          ]
            \addplot[color=ThesisPrimary, line width=1.4pt, smooth]
              coordinates {(1,68.5) (4,84.7) (8,94.1) (12,97.8) (16,99.1) (20,99.7)};
            \addlegendentry{Méthode proposée}
            \addplot[color=ThesisAccent, line width=1.1pt, dashed, smooth]
              coordinates {(1,56.7) (4,67.1) (8,75.1) (12,79.3) (16,81.5) (20,82.7)};
            \addlegendentry{Référence}
          \end{axis}
        \end{tikzpicture}
      \end{cadreFigure}
    \end{column}
    \begin{column}{0.46\textwidth}
      \begin{slidecard}{Critère d'optimisation}
        Minimisation régularisée à chaque étape~\autocite{tantau2023tikz} :
        \[
          \vect{\theta}_{k+1} = \argmin_{\vect{\theta} \in \setR^d}
          \bigl\{ \mathcal{L}(\vect{\theta}) + \tfrac{\lambda}{2}\norm{\vect{\theta}}_2^2 \bigr\}
        \]
      \end{slidecard}
      \vspace*{0.15cm}
      \begin{callout}[Apparition progressive]
        Les overlays natifs Beamer (\code{\textbackslash pause}, \code{<2->}) guident le regard sans surcharger l'écran.
      \end{callout}
    \end{column}
  \end{columns}
\end{frame}

\begin{frame}[fragile]{Extrait de code lisible et isolé du contenu}
  \framesubtitle{Coloration syntaxique cohérente et police monospace dédiée}

  \begin{columns}[T, totalwidth=\textwidth]
    \begin{column}{0.62\textwidth}
\begin{lstlisting}[language=Python]
def build_deck(source: str, ratio: str = "169") -> str:
  """Compile et audite un support éditorial LuaLaTeX."""
  options = ["french", f"ratio={ratio}", "theme=light"]
  return f"\\documentclass[{', '.join(options)}]{{ClasseGuillaumeSlides}}"
\end{lstlisting}
    \end{column}
    \begin{column}{0.35\textwidth}
      \begin{alertcard}{Règle \code{fragile}}
        Ajouter impérativement \code{[fragile]} sur toute \code{frame} contenant un bloc \code{lstlisting} ou \code{verbatim}.
      \end{alertcard}
    \end{column}
  \end{columns}

  \takeaway{Les extraits longs sont placés dans \code{listings/*.py} et importés via \code{\textbackslash lstinputlisting}.}
\end{frame}
